\scanpage{52}
\mspage{11}
(il y a \unsure{princip\textsuperscript{lt}|principalement|manifestement} \unsure{2 cas|le cas|4 cas} à distinguer, ceux des $r=r_{ij}$ et $r=r_{ijh}$ \ldots) dans le cas standard,
\[
\alpha=\alpha^{ij}=(\zeta_i,\zeta'_j,(\unsure{b'_{ih}|\zeta'_{ih}|\xi_{ik}})_{h\ne i,j})\quad\text{et}\quad
\alpha=\alpha^{i,j,h}=(\zeta_i,\zeta_j,\zeta_h,\zeta_{\ell m},\zeta_{mn},\zeta_{\unsure{np|n\ell}})
\]
sont des bases. \deleted{\unsure{si|i.e.}} \inserted{Soit} \deleted{si} \deleted{$\alpha=(\alpha_1,\ldots\alpha_6)$} $\alpha'=\{\alpha'_1,\ldots,\alpha'_6\}$,

\deleted{$\Rightarrow\ \forall\ 1\le i<j\le6,\ \exists!\ \unsure{\delta}\in\alpha_{ij}\in\Delta$ lié à \unsure{$\zeta_i$} $\alpha_i,\alpha'_j$}

Si $1\le i<j\le6$, \deleted{$\unsure{\delta}$}$\alpha_i$ et $\alpha'_j$ sont liés, donc $\exists!\ \delta\in\Delta_0$ lié à $\alpha_i$ et $\alpha'_j$, soit $\alpha_{ij}$. \deleted{Les $\alpha_{ij}$} Tout $\delta\in\Delta_0$ distinct des $\alpha_i,\alpha'_i$ est-il de la forme $\alpha_{ij}$ i.e. est-il lié à un des $\alpha_i$ et à un des $\alpha'_j$ ? \deleted{S'il n'était lié à aucun des \ldots}

\deleted{Vérifications fastidieuses \ldots} Meilleur procédé :

Comme pour tte racine $r$ on a $rr=-2$, l'endomorphisme (``symétrie \unsure{parr à|par à|par} $r$'')
\[
s_r\deleted{\unsure{(|:}}x\longmapsto x+(x.r)r\qquad\widehat E\longrightarrow\widehat E
\]
[dont l'ens des p\textsuperscript{ts} fixes est $r^\perp$, et qui transforme $r$ en $-r$] est un autom de \deleted{orthogonal} \inserted{$(\widehat E,\widehat\varphi\ \text{forme fond.})$} qui évidemment respecte aussi $\lambda$ (car $r\in E=\lambda^\perp$). Montrons qu'elle applique $\Delta$ dans $\Delta$. [Il suffit de regarder le cas standard, et de regarder alors les trois cas $r=r_0$, $r=r_{ij}$, $r=r_{ijh}$.]

\deleted{\unsure{$r_0$}} Notons d'abord, de façon générale, que si $r$ est associé à $(\alpha_1,\ldots\alpha_6),(\alpha'_1,\ldots\alpha'_6)$ ($r=\alpha'_i-\alpha_i$ pour $1\le i\le\unsure{r|6}$) on a
\[
\boxed{s_r\alpha_i=\alpha'_i,\quad s_r\alpha'_i=\alpha_i}
\]
Car $\left\{\begin{aligned}r.\alpha_i&=1,\ r\alpha'_i=-1\end{aligned}\right.$

\deleted{Il reste} à voir l'effet de $r$ sur les autres éléments de $\Delta$ --- et de montrer que
\[
\boxed{s_r\delta=\delta\quad\text{si }\delta\notin\{(\alpha_i),(\alpha'_i)\}}
\]
% archive sheet number: 51
